% ==============================================================================
% CHƯƠNG 2: TÍNH TOÁN LÝ THUYẾT (CÂU a, b, c)
% ==============================================================================
\section{Tính toán lý thuyết (Câu a, b, c)}

\subsection{Câu a: Khai triển chuỗi Fourier lượng giác cho nguồn áp \texorpdfstring{$e(t)$}{e(t)}}

Chuỗi Fourier lượng giác tổng quát của tín hiệu tuần hoàn $e(t)$ có chu kỳ $T = 1/f_0$ và tần số góc cơ bản $\omega_0 = 2\pi f_0$ có dạng:
\begin{equation}
    e(t) = a_0 + \sum_{n=1}^{\infty} \Big[ a_n \cos(n\omega_0 t) + b_n \sin(n\omega_0 t) \Big].
    \label{eq:fourier_general}
\end{equation}
Đặt biến góc pha $\theta = \omega_0 t \in [0, 2\pi]$, ta có $\dd t = \dd\theta / \omega_0$ và $e(\theta) = E_m \sin\theta$ với $\theta \in [0, \pi)$, $e(\theta) = 0$ với $\theta \in [\pi, 2\pi)$.

\subsubsection*{1. Tính thành phần một chiều (DC) $a_0$}
\begin{equation}
    a_0 = \frac{1}{T}\int_{0}^{T} e(t)\dd t
        = \frac{1}{2\pi}\int_{0}^{\pi} E_m \sin\theta \dd\theta
        = \frac{E_m}{2\pi}\Big[-\cos\theta\Big]_{0}^{\pi}
        = \frac{E_m}{\pi}.
\end{equation}

\subsubsection*{2. Tính hệ số $a_1, b_1$ ứng với sóng hài cơ bản ($n = 1$)}
\begin{align}
    a_1 &= \frac{2}{T}\int_{0}^{T/2} E_m \sin(\omega_0 t)\cos(\omega_0 t)\dd t
         = \frac{E_m}{2\pi}\int_{0}^{\pi} \sin(2\theta)\dd\theta
         = \frac{E_m}{2\pi}\left[-\frac{\cos(2\theta)}{2}\right]_{0}^{\pi} = 0, \\
    b_1 &= \frac{2}{T}\int_{0}^{T/2} E_m \sin^2(\omega_0 t)\dd t
         = \frac{E_m}{\pi}\int_{0}^{\pi} \frac{1 - \cos(2\theta)}{2}\dd\theta
         = \frac{E_m}{2\pi}\left[\theta - \frac{\sin(2\theta)}{2}\right]_{0}^{\pi}
         = \frac{E_m}{2}.
\end{align}

\subsubsection*{3. Tính hệ số $a_n, b_n$ với $n \ge 2$}
Sử dụng công thức tích thành tổng lượng giác:
\begin{align}
    b_n &= \frac{E_m}{\pi}\int_{0}^{\pi} \sin\theta \sin(n\theta)\dd\theta
         = \frac{E_m}{2\pi}\int_{0}^{\pi} \Big[\cos\big((n-1)\theta\big) - \cos\big((n+1)\theta\big)\Big]\dd\theta \nonumber\\
        &= \frac{E_m}{2\pi}\left[\frac{\sin\big((n-1)\theta\big)}{n-1} - \frac{\sin\big((n+1)\theta\big)}{n+1}\right]_{0}^{\pi} = 0 \qquad (\forall n \ge 2).
\end{align}
Tương tự với $a_n$:
\begin{align}
    a_n &= \frac{E_m}{\pi}\int_{0}^{\pi} \sin\theta \cos(n\theta)\dd\theta
         = \frac{E_m}{2\pi}\int_{0}^{\pi} \Big[\sin\big((n+1)\theta\big) - \sin\big((n-1)\theta\big)\Big]\dd\theta \nonumber\\
        &= \frac{E_m}{2\pi}\left[-\frac{\cos\big((n+1)\theta\big)}{n+1} + \frac{\cos\big((n-1)\theta\big)}{n-1}\right]_{0}^{\pi} \nonumber\\
        &= \frac{E_m}{2\pi}\left[\frac{1 - \cos\big((n+1)\pi\big)}{n+1} - \frac{1 - \cos\big((n-1)\pi\big)}{n-1}\right].
\end{align}
Xét hai trường hợp theo tính chẵn lẻ của $n \ge 2$:
\begin{itemize}
    \item Nếu $n$ \textbf{lẻ} ($n = 3, 5, 7, \dots$): $n+1$ và $n-1$ đều chẵn nên $\cos((n\pm 1)\pi) = +1 \implies a_n = 0$.
    \item Nếu $n$ \textbf{chẵn} ($n = 2k,\; k = 1, 2, 3, \dots$): $n+1$ và $n-1$ đều lẻ nên $\cos((n\pm 1)\pi) = -1$, do đó:
    \begin{equation}
        a_n = \frac{E_m}{2\pi}\left(\frac{2}{n+1} - \frac{2}{n-1}\right)
            = -\frac{2E_m}{\pi(n^2 - 1)}
            = -\frac{2E_m}{\pi(4k^2 - 1)}.
    \end{equation}
\end{itemize}

\subsubsection*{4. Kết quả chuỗi Fourier lượng giác của nguồn $e(t)$}
\begin{equation}
    e(t) = \frac{E_m}{\pi} + \frac{E_m}{2}\sin(\omega_0 t)
           - \frac{2E_m}{\pi}\sum_{k=1}^{\infty} \frac{\cos(2k\omega_0 t)}{4k^2 - 1} \quad (\text{V}).
    \label{eq:e_fourier_sym}
\end{equation}
Thay $E_m = 19\text{ V}$ và $\omega_0 = 160\pi\text{ rad/s}$:
\begin{equation}
    e(t) = \frac{19}{\pi} + 9{,}5\sin(160\pi t)
           - \frac{38}{\pi}\sum_{k=1}^{\infty} \frac{\cos(320k\pi t)}{4k^2 - 1} \quad (\text{V}),
\end{equation}
hay viết tường minh các thành phần đầu tiên:
\begin{align}
    e(t) &= 6{,}0479 + 9{,}5000\sin(160\pi t) - 4{,}0319\cos(320\pi t) - 0{,}8064\cos(640\pi t) \nonumber\\
         &\quad - 0{,}3456\cos(960\pi t) - 0{,}1920\cos(1280\pi t) - \cdots \quad (\text{V}).
\end{align}
Để tiện cho phép tính miền phức (phasor) và mô phỏng nguồn \texttt{SINE} trên LTspice, có thể quy đổi toàn bộ về biên độ dương dựa trên đồng nhất thức $-\cos\alpha = \cos(\alpha + 180^\circ) = \sin(\alpha - 90^\circ)$ và $\sin\alpha = \cos(\alpha - 90^\circ)$.

Hình~\ref{fig:source_recon} minh họa nguồn chỉnh lưu bán kỳ thực tế $e(t)$ cùng dạng sóng xấp xỉ khi lấy thành phần DC kết hợp 3 sóng hài ($n=1,2,4$) và 5 sóng hài ($n=1,2,4,6,8$) đầu tiên.

\begin{figure}[!htbp]
    \centering
    \begin{tikzpicture}
        \begin{axis}[
            width=14.6cm, height=5.8cm,
            xlabel={Thời gian $t$ (ms)},
            ylabel={Điện áp nguồn $e(t)$ (V)},
            xmin=0, xmax=25,
            ymin=-2.5, ymax=21.5,
            xtick={0,3.125,6.25,9.375,12.5,15.625,18.75,21.875,25},
            xticklabels={0,3.125,6.25,9.375,12.5,15.625,18.75,21.875,25},
            grid=both,
            grid style={gray!25},
            legend pos=north east,
            legend style={font=\small, fill=white, fill opacity=0.9}
        ]
            \addplot[black, line width=1.3pt] table[x=t_ms, y=e_real] {data/waveforms_2T.dat};
            \addlegendentry{Nguồn thực $e(t)$}
            \addplot[accentred, dashed, line width=1.0pt] table[x=t_ms, y=e_3h] {data/waveforms_2T.dat};
            \addlegendentry{Xấp xỉ DC + 3 hài ($n=1,2,4$)}
            \addplot[hcmutblue, dotted, line width=1.2pt] table[x=t_ms, y=e_5h] {data/waveforms_2T.dat};
            \addlegendentry{Xấp xỉ DC + 5 hài ($n=1,2,4,6,8$)}
        \end{axis}
    \end{tikzpicture}
    \caption{Dạng sóng nguồn áp chỉnh lưu bán kỳ $e(t)$ và chuỗi Fourier xấp xỉ trong 2 chu kỳ ($2T = 25\text{ ms}$).}
    \label{fig:source_recon}
\end{figure}

% ==============================================================================
\subsection{Câu b: Biểu thức tổng quát cho điện áp \texorpdfstring{$u_0(t)$}{u0(t)} dưới dạng chuỗi vô hạn}

Áp dụng phương pháp phức hóa và nguyên lý xếp chồng cho mạch tuyến tính bất biến theo thời gian ở chế độ xác lập điều hòa.

\subsubsection*{1. Xét thành phần một chiều ($n = 0$, $\omega = 0$)}
Ở chế độ xác lập một chiều, tụ điện $C$ có tổng trở $Z_C(0) = \infty$ (hở mạch) và cuộn cảm $L$ có tổng trở $Z_L(0) = 0$ (ngắn mạch), do đó dòng điện một chiều bằng $0$ và toàn bộ điện áp DC $a_0 = E_m/\pi$ đặt lên tụ $C$:
\begin{equation}
    H(0) = 0 \implies U_{0,0} = 0\text{ V}.
\end{equation}

\subsubsection*{2. Xét thành phần xoay chiều bậc $n \ge 1$ (tần số góc $\omega_n = n\omega_0$)}
Tại tần số góc $\omega_n = n\omega_0$, tổng trở phức của tụ $C$ và nhánh song song $(R \parallel L)$ lần lượt là:
\begin{equation}
    Z_C(\jj n\omega_0) = \frac{1}{\jj n\omega_0 C}, \qquad
    Z_{RL}(\jj n\omega_0) = \frac{R \cdot \jj n\omega_0 L}{R + \jj n\omega_0 L}.
\end{equation}
Theo quy tắc phân áp phức, hàm truyền điện áp $H(\jj n\omega_0) = \dot{U}_{0,n}/\dot{E}_n$ là:
\begin{align}
    H(\jj n\omega_0)
    &= \frac{Z_{RL}(\jj n\omega_0)}{Z_C(\jj n\omega_0) + Z_{RL}(\jj n\omega_0)}
     = \frac{\dfrac{\jj n\omega_0 RL}{R + \jj n\omega_0 L}}{\dfrac{1}{\jj n\omega_0 C} + \dfrac{\jj n\omega_0 RL}{R + \jj n\omega_0 L}} \nonumber\\
    &= \frac{(\jj n\omega_0)^2 RLC}{R + \jj n\omega_0 L + (\jj n\omega_0)^2 RLC}
     = \frac{-(n\omega_0)^2 LC}{\Big[1 - (n\omega_0)^2 LC\Big] + \jj\,\dfrac{n\omega_0 L}{R}}.
    \label{eq:transfer_H}
\end{align}

Từ \eqref{eq:transfer_H}, mô-đun và góc pha của hàm truyền $H(\jj n\omega_0)$ tại bậc hài $n \ge 1$ được xác định chính xác bằng:
\begin{align}
    |H(\jj n\omega_0)| &= \frac{(n\omega_0)^2 LC}{\sqrt{\Big[1 - (n\omega_0)^2 LC\Big]^2 + \left(\dfrac{n\omega_0 L}{R}\right)^2}}, \label{eq:H_mag}\\
    \varphi_{H,n} &= \arg H(\jj n\omega_0) = 180^\circ - \operatorname{atan2}\!\left(\frac{n\omega_0 L}{R},\; 1 - (n\omega_0)^2 LC\right) \in (0^\circ,\, 180^\circ). \label{eq:H_phase}
\end{align}
\textit{(Kiểm chứng giới hạn vật lý: khi $n\omega_0 \to 0^+$ thì $|H| \to 0,\; \varphi_H \to 180^\circ$; tại tần số cộng hưởng $\omega_r = 1/\sqrt{LC}$ thì $\varphi_H = 90^\circ$; khi $n\omega_0 \to \infty$ thì $|H| \to 1,\; \varphi_H \to 0^\circ$ -- hoàn toàn đúng bản chất bộ lọc thông cao bậc 2.)}

\subsubsection*{3. Biểu thức tổng quát của đáp ứng xác lập $u_0(t)$}
Áp dụng nguyên lý xếp chồng cho từng thành phần của chuỗi nguồn \eqref{eq:e_fourier_sym}, ta thu được đáp ứng điện áp xác lập dưới dạng chuỗi vô hạn:
\begin{equation}
    u_0(t) = \frac{E_m}{2}\,|H(\jj\omega_0)|\sin(\omega_0 t + \varphi_{H,1})
           - \frac{2E_m}{\pi}\sum_{k=1}^{\infty} \frac{|H(\jj 2k\omega_0)|}{4k^2 - 1}\cos(2k\omega_0 t + \varphi_{H,2k})
           \quad (\text{V}).
    \label{eq:u0_inf}
\end{equation}
Viết gọn toàn bộ về dạng chuẩn biên độ dương theo hàm $\cos$:
\begin{equation}
    u_0(t) = U_{0,1}\cos(\omega_0 t + \psi_1) + \sum_{k=1}^{\infty} U_{0,2k}\cos(2k\omega_0 t + \psi_{2k}) \quad (\text{V}),
    \label{eq:u0_inf_cos}
\end{equation}
trong đó:
\begin{align}
    U_{0,1} &= \frac{E_m}{2}|H(\jj\omega_0)|, & \psi_1 &= \varphi_{H,1} - 90^\circ, \\
    U_{0,2k} &= \frac{2E_m}{\pi(4k^2 - 1)}|H(\jj 2k\omega_0)|, & \psi_{2k} &= \varphi_{H,2k} - 180^\circ \quad (k = 1, 2, 3, \dots).
\end{align}

Hình~\ref{fig:bode_spectrum} trình bày đặc tuyến biên độ $|H(f)|$ của mạch cùng phổ biên độ nguồn $|E_n|$ và ngõ ra $U_{0,n}$. Có thể thấy rõ vùng cộng hưởng quanh $f_r \approx 317{,}04\text{ Hz}$ nâng mạnh đóng góp của các sóng hài bậc $4$, $6$, $8$, đồng thời làm cho biên độ hài bậc 2 ($U_{0,2} = 1{,}1423\text{ V}$) và hài bậc 4 ($U_{0,4} = 0{,}8170\text{ V}$) ở ngõ ra thậm chí lớn hơn cả hài cơ bản bậc 1 ($U_{0,1} = 0{,}6239\text{ V}$).

\begin{figure}[!htbp]
    \centering
    \begin{tikzpicture}
        \begin{axis}[
            name=bodeplot,
            width=6.4cm, height=5.4cm,
            xlabel={Tần số $f$ (Hz)},
            ylabel={Độ lớn hàm truyền $|H(\jj 2\pi f)|$},
            xmin=0, xmax=1000,
            ymin=0, ymax=1.32,
            grid=both,
            grid style={gray!25},
            title={\small (a) Đặc tuyến biên độ $|H(f)|$},
            legend pos=south east,
            legend style={font=\scriptsize}
        ]
            \addplot[hcmutblue, line width=1.2pt] table[x=f_Hz, y=H_mag] {data/bode.dat};
            \addlegendentry{$|H(f)|$}
            \addplot[accentred, only marks, mark=*, mark size=2pt] table[x=f_Hz, y=H_mag] {data/spectrum.dat};
            \addlegendentry{Các hài $n=0,1,2,4,\dots$}
            \draw[accentgreen, dashed, line width=0.8pt] (axis cs:317.04,0) -- (axis cs:317.04,1.32);
            \node[accentgreen, anchor=south west, font=\scriptsize] at (axis cs:325,0.06) {$f_r \approx 317\text{ Hz}$};
        \end{axis}

        \begin{axis}[
            at={($(bodeplot.east)+(1.7cm,0)$)},
            anchor=west,
            width=6.4cm, height=5.4cm,
            ybar,
            bar width=7pt,
            enlarge x limits=0.14,
            xlabel={Bậc sóng hài $n$},
            ylabel={Biên độ đỉnh $U_{0,n}$ (V)},
            symbolic x coords={0,1,2,4,6,8,10,12},
            xtick=data,
            ymin=0, ymax=1.38,
            ymajorgrids=true,
            grid style={gray!25},
            title={\small (b) Phổ biên độ ngõ ra $U_{0,n}$},
            nodes near coords,
            every node near coord/.append style={
                font=\tiny,
                rotate=45,
                anchor=west,
                /pgf/number format/fixed,
                /pgf/number format/precision=2
            },
            nodes near coords align={vertical}
        ]
            \addplot[fill=hcmutblue!75, draw=hcmutblue] table[x=n, y=U0_amp] {data/spectrum.dat};
        \end{axis}
    \end{tikzpicture}
    \caption{(a) Đặc tuyến biên độ bộ lọc thông cao bậc 2 $|H(f)|$ với $f_r \approx 317{,}04\text{ Hz}$ và (b) phổ biên độ điện áp ngõ ra $U_{0,n}$ tại các sóng hài khác $0$ đầu tiên.}
    \label{fig:bode_spectrum}
\end{figure}

% ==============================================================================
\subsection{Câu c: Xấp xỉ \texorpdfstring{$u_0(t)$}{u0(t)} với 3 sóng hài và 5 sóng hài khác 0 đầu tiên}

Do thành phần DC ($n=0$) bị tụ $C$ chặn hoàn toàn ($U_{0,0} = 0$) và các sóng hài lẻ bậc $n \ge 3$ có hệ số nguồn bằng $0$, nên:
\begin{itemize}
    \item \textbf{3 sóng hài khác $0$ đầu tiên} của $u_0(t)$ gồm $n \in \{1, 2, 4\}$ (ứng với tần số $80\text{ Hz}$, $160\text{ Hz}$, $320\text{ Hz}$).
    \item \textbf{5 sóng hài khác $0$ đầu tiên} của $u_0(t)$ gồm $n \in \{1, 2, 4, 6, 8\}$ (ứng với tần số $80\text{ Hz}$, $160\text{ Hz}$, $320\text{ Hz}$, $480\text{ Hz}$, $640\text{ Hz}$).
\end{itemize}

Thay số $E_m = 19\text{ V}$, $f_0 = 80\text{ Hz}$, $R = 28\ \Omega$, $L = 14\text{ mH}$, $C = 18\ \mu\text{F}$ vào \eqref{eq:H_mag}--\eqref{eq:u0_inf_cos}, ta lập được bảng thông số từng hài như trong Bảng~\ref{tab:harmonics}.

\begin{table}[!htbp]
    \centering
    \caption{Tổng hợp thông số các sóng hài khác $0$ đầu tiên của nguồn $e(t)$ và đáp ứng $u_0(t)$ (Đề~5).}
    \label{tab:harmonics}
    \vspace{4pt}
    \small
    \begin{tabular}{|c|r|r|r|r|r|r|r|}
        \toprule
        \textbf{Bậc $n$} & \textbf{$f_n$ (Hz)} & \textbf{Biên độ $E_n$ (V)} & \textbf{$|H(\jj n\omega_0)|$} & \textbf{$\varphi_{H,n}$ ($^\circ$)} & \textbf{$U_{0,n}$ (V)} & \textbf{$\psi_n$ ($\cos$, $^\circ$)} & \textbf{$\psi_n$ ($\sin$, $^\circ$)} \\
        \midrule
        $0$ (DC) & $0$   & $6{,}0479$ & $0{,}0000$ & --          & $0{,}0000$ & --          & -- \\
        \hline
        $1$      & $80$  & $9{,}5000$ & $0{,}0657$ & $+164{,}97$ & $0{,}6239$ & $+74{,}97$  & $+164{,}97$ \\
        $2$      & $160$ & $4{,}0319$ & $0{,}2833$ & $+146{,}00$ & $1{,}1423$ & $-34{,}00$  & $+56{,}00$  \\
        $4$      & $320$ & $0{,}8064$ & $1{,}0132$ & $+88{,}93$  & $0{,}8170$ & $-91{,}07$  & $-1{,}07$   \\
        \hline
        $6$      & $480$ & $0{,}3456$ & $1{,}1542$ & $+49{,}41$  & $0{,}3989$ & $-130{,}59$ & $-40{,}59$  \\
        $8$      & $640$ & $0{,}1920$ & $1{,}1091$ & $+33{,}18$  & $0{,}2130$ & $-146{,}82$ & $-56{,}82$  \\
        \hline
        \multicolumn{8}{|l|}{\textit{Tham khảo thêm hai hài kế tiếp (đánh giá đuôi chuỗi):}} \\
        $10$     & $800$ & $0{,}1222$ & $1{,}0744$ & $+25{,}09$  & $0{,}1313$ & $-154{,}91$ & $-64{,}91$  \\
        $12$     & $960$ & $0{,}0846$ & $1{,}0529$ & $+20{,}26$  & $0{,}0891$ & $-159{,}74$ & $-69{,}74$  \\
        \bottomrule
    \end{tabular}
\end{table}

\subsubsection*{1. Biểu thức xấp xỉ $u_0(t)$ với 3 sóng hài khác $0$ đầu tiên ($n = 1, 2, 4$)}
Viết theo dạng chuẩn hàm $\cos$ (biên độ dương):
\begin{align}
    u_{0,3\text{h}}(t)
    &= 0{,}6239\cos(160\pi t + 74{,}97^\circ)
     + 1{,}1423\cos(320\pi t - 34{,}00^\circ) \nonumber\\
    &\quad + 0{,}8170\cos(640\pi t - 91{,}07^\circ) \quad (\text{V}),
    \label{eq:u0_3h}
\end{align}
(tương đương theo dạng \eqref{eq:u0_inf}: $0{,}6239\sin(160\pi t + 164{,}97^\circ) - 1{,}1423\cos(320\pi t + 146{,}00^\circ) - 0{,}8170\cos(640\pi t + 88{,}93^\circ)\text{ V}$).

Giá trị hiệu dụng theo định lý Parseval:
\begin{equation}
    U_{0,3\text{h},\RMS} = \sqrt{\frac{1}{2}\left(0{,}6239^2 + 1{,}1423^2 + 0{,}8170^2\right)} = 1{,}0866\text{ V}.
\end{equation}

\subsubsection*{2. Biểu thức xấp xỉ $u_0(t)$ với 5 sóng hài khác $0$ đầu tiên ($n = 1, 2, 4, 6, 8$)}
\begin{align}
    u_{0,5\text{h}}(t)
    &= 0{,}6239\cos(160\pi t + 74{,}97^\circ)
     + 1{,}1423\cos(320\pi t - 34{,}00^\circ) \nonumber\\
    &\quad + 0{,}8170\cos(640\pi t - 91{,}07^\circ)
     + 0{,}3989\cos(960\pi t - 130{,}59^\circ) \nonumber\\
    &\quad + 0{,}2130\cos(1280\pi t - 146{,}82^\circ) \quad (\text{V}).
    \label{eq:u0_5h}
\end{align}

Giá trị hiệu dụng theo định lý Parseval:
\begin{equation}
    U_{0,5\text{h},\RMS} = \sqrt{\frac{1}{2}\left(0{,}6239^2 + 1{,}1423^2 + 0{,}8170^2 + 0{,}3989^2 + 0{,}2130^2\right)} = 1{,}1327\text{ V},
\end{equation}
so với giá trị hiệu dụng lý thuyết chính xác (tính với chuỗi vô hạn) là $U_{0,\infty,\RMS} = 1{,}1408\text{ V}$.
